\documentclass[11pt]{article}
\usepackage[margin=1.2in]{geometry}
\usepackage{amsmath,amssymb,amsthm}
\usepackage{booktabs}
\usepackage{xcolor}
\definecolor{linkblue}{RGB}{25,60,120}
\usepackage[colorlinks=true,linkcolor=linkblue,citecolor=linkblue,urlcolor=linkblue]{hyperref}
\usepackage{microtype}

\newtheorem{theorem}{Theorem}
\newtheorem{lemma}[theorem]{Lemma}
\newtheorem{proposition}[theorem]{Proposition}
\theoremstyle{remark}
\newtheorem{remark}[theorem]{Remark}

\newcommand{\I}{\mathrel{;}}
\DeclareMathOperator{\Ent}{H}

\title{An explicit family of classical states refuting a conjecture of Linden and Winter}
\author{Drew Stone\thanks{\texttt{drewstone329@gmail.com}. Code and machine-checkable
certificates: \url{https://github.com/drewstone/lw6-counterexample}.}}
\date{August 19, 2026}

\begin{document}
\maketitle

\begin{abstract}
Linden and Winter (\emph{Commun.\ Math.\ Phys.}\ 259:129--138, 2005) proved that any
four-party quantum state with $I(A\I C|B)=I(C\I B|A)=I(A\I B|D)=0$ satisfies
$I(C\I D)\ge I(C\I AB)$, and conjectured (their Conjecture~6) that this constrained
inequality is stable: that positive constants $k_1,k_2,k_3$ exist with
$k_1 I(A\I C|B)+k_2 I(C\I B|A)+k_3 I(A\I B|D)+I(C\I D)-I(C\I AB)\ge 0$
for all quadripartite states. We exhibit an explicit one-parameter family of
\emph{classical} probability distributions on four bits --- six atoms, integer weights in
closed form --- on which $I(A\I C|B)$ and $I(A\I B|D)$ vanish identically while the ratio
$\bigl[I(C\I AB)-I(C\I D)\bigr]/I(C\I B|A)$ grows without bound. Consequently no choice of
positive constants satisfies the conjectured inequality, even for classical states: the
conjecture is false. The vanishing of the two conditional mutual informations is an exact
combinatorial identity; the divergence of the ratio follows from an elementary asymptotic
analysis, and is independently certified by directed-rounding interval arithmetic, which
also yields explicit refuting states for every $(k_1,k_2,k_3)$ with
$k_2\le 10^{24}$. All certificates re-verify from a single dependency-free Python script.
\end{abstract}

\section{Introduction}

Which linear inequalities constrain the von Neumann entropies of the reduced states of a
multipartite quantum state? For up to three parties the answer is classical: strong
subadditivity and weak monotonicity~\cite{LiebRuskai1973} generate all inequalities, and
Pippenger~\cite{Pippenger2003} asked whether anything new appears for four or more
parties. No unconstrained inequality beyond the Lieb--Ruskai family has been found since
1973.

Linden and Winter~\cite{LindenWinter2005} made the first advance past this point,
proving a new \emph{constrained} inequality:

\begin{quote}
\textbf{Theorem (Linden--Winter).} If a four-party state $\rho_{ABCD}$ satisfies
$I(A\I C|B)=I(C\I B|A)=I(A\I B|D)=0$, then $I(C\I D)\ge I(C\I AB)$.
\end{quote}

\noindent
They closed the paper by conjecturing that the constraint can be relaxed linearly:

\begin{quote}
\textbf{Conjecture 6 (Linden--Winter).} There exist positive real constants
$k_1,k_2,k_3$ such that for all quadripartite quantum states,
\begin{equation}\label{eq:lw6}
k_1\, I(A\I C|B) + k_2\, I(C\I B|A) + k_3\, I(A\I B|D) + I(C\I D) - I(C\I AB)\;\ge\;0 .
\end{equation}
\end{quote}

\noindent
A proof would have produced the first new unconstrained inequality for the von Neumann
entropy in over five decades. We show the conjecture fails, already for classical states.

\begin{theorem}\label{thm:main}
No positive constants $k_1,k_2,k_3$ satisfy~\eqref{eq:lw6}. Specifically, for the family
$W_L$ of classical distributions in Definition~\ref{def:family}: $I(A\I C|B)$ and
$I(A\I B|D)$ vanish identically on the family, and
$R_L := \bigl[I(C\I AB)-I(C\I D)\bigr]/I(C\I B|A) \longrightarrow \infty$
as $L\to\infty$. Hence for any positive $(k_1,k_2,k_3)$, every $L$ with $R_L > k_2$ gives
$f_k(W_L)<0$ in~\eqref{eq:lw6}.
\end{theorem}

Classical distributions embed as diagonal quantum states, so a classical counterexample
refutes the quantum conjecture. The constrained theorem itself is untouched: on the
family, two of its three hypotheses hold exactly but the third,
$I(C\I B|A)=0$, fails --- by an amount that shrinks too fast, which is precisely what the
conjecture forbade. The result says the Linden--Winter inequality is not linearly stable:
the deficit $I(C\I AB)-I(C\I D)$ can exceed any fixed multiple of the constraint
violations.

\paragraph{Provenance.} The family was found by an automated research system conducting a
systematic search over supports and weight patterns, and the numerical certificates below
were produced and cross-checked by four independently written implementations (no shared
code), including the dependency-free verifier published with this note. The mathematics
presented here --- the identities of Lemma~\ref{lem:zeros} and the asymptotics of
Proposition~\ref{prop:asym} --- is human-checkable by hand.

\section{The family}

\begin{samepage}
\begin{proposition}[Definition]\label{def:family}
For an integer $L\ge 2$ let $W_L$ be the probability distribution on
$\{0,1\}^4$ (bits ordered $A,B,C,D$) supported on six atoms, with weights proportional
to:
\begin{center}
\begin{tabular}{cc}
\toprule
atom $(A,B,C,D)$ & weight \\
\midrule
$(0,0,0,0)$ & $2\cdot 10^{\,4L+3}$ \\
$(0,0,1,0)$ & $27\cdot 10^{\,4L+3}-27\cdot 10^{\,3L+3}$ \\
$(1,1,0,1)$ & $2\cdot 10^{\,4L}$ \\
$(1,1,1,1)$ & $27\cdot 10^{\,4L}$ \\
$(0,1,0,1)$ & $2\cdot 10^{\,L+3}$ \\
$(0,1,1,0)$ & $27\cdot 10^{\,L+3}$ \\
\bottomrule
\end{tabular}
\end{center}
Equivalently, with $\delta=10^{-L}$, the weights are proportional to
$2,\;27(1-\delta),\;\tfrac{2}{1000},\;\tfrac{27}{1000},\;2\delta^{3},\;27\delta^{3}$
on the same atoms.
\end{proposition}
\end{samepage}

\begin{lemma}[Two exact zeros]\label{lem:zeros}
For every $L$: $I(A\I B|D)=0$ and $I(A\I C|B)=0$ on $W_L$, exactly.
\end{lemma}

\begin{proof}
$I(A\I B|D)$: on the three atoms with $D=0$ the bit $A$ is constant ($=0$), and on the
three atoms with $D=1$ the bit $B$ is constant ($=1$). Conditioned on either value of
$D$, one of the two arguments of the mutual information is deterministic, so each
conditional term vanishes.

$I(A\I C|B)$: on the two atoms with $B=0$ the bit $A$ is constant ($=0$), so that block
contributes zero. On the $B=1$ block the joint distribution of $(A,C)$ has the $2\times2$
weight matrix
$\bigl(\begin{smallmatrix} w_{0101} & w_{0110}\\ w_{1101} & w_{1111}\end{smallmatrix}\bigr)$,
and $A$ and $C$ are conditionally independent iff its determinant vanishes:
$w_{1101}\,w_{0110} = w_{1111}\,w_{0101}$. Both products equal
$54\cdot 10^{\,5L+3}$.
\end{proof}

Because of Lemma~\ref{lem:zeros}, on the family the left side of~\eqref{eq:lw6} collapses
to
\begin{equation}\label{eq:collapse}
f_k(W_L) \;=\; k_2\, y_L + e_L,
\qquad y_L := I(C\I B|A),\qquad e_L := I(C\I D)-I(C\I AB),
\end{equation}
for \emph{every} $(k_1,k_3)$: those two coefficients multiply exact zeros.

\section{Why the ratio diverges}

Write $\delta=10^{-L}$ and use the normalized form of Definition~\ref{def:family}.
Two structural observations drive everything.

\begin{proposition}\label{prop:asym}
As $\delta\to 0$: $y_L=\Theta(\delta^{5})$ while $-e_L=\Theta(\delta^{4})$, with $e_L<0$
for all small $\delta$. Consequently $R_L=-e_L/y_L=\Theta(1/\delta)\to\infty$.
\end{proposition}

\begin{proof}[Proof sketch, elementary and self-contained]
\emph{(i) $y_L$ needs two rare events at once.} Only the $A{=}0$ block contributes to
$I(C\I B|A)$ (given $A{=}1$, $B\equiv 1$). Within $A{=}0$, the branch $B{=}1$ carries
probability $\Theta(\delta^{3})$ (the two $\delta^3$-atoms), and conditioned on the two
branches the distributions of $C$ are $(2,27)/29$ versus
$\bigl(2,\,27(1-\delta)\bigr)/(29-27\delta)$ --- total-variation distance
$\Theta(\delta)$. A mutual information with a $\Theta(\delta^3)$-rare branch whose
conditional laws differ by $\Theta(\delta)$ is
$\Theta(\delta^{3})\cdot\Theta(\delta^{2}) = \Theta(\delta^{5})$, by the local
$\chi^2$-expansion of Kullback--Leibler divergence.

\emph{(ii) $e_L$ vanishes exactly when the $\delta^3$-atoms are removed.} Set the two
$\delta^{3}$-weights to zero. Then only four atoms remain, on which the partition of the
sample space induced by $D$ \emph{coincides} with the partition induced by $(A,B)$: the
$D{=}0$ atoms are exactly the $AB{=}00$ atoms and the $D{=}1$ atoms are exactly the
$AB{=}11$ atoms. Hence $I(C\I D)=I(C\I AB)$ identically in $\delta$ on the truncated
family, i.e.\ $e\equiv 0$. Restoring the two atoms perturbs each side of
$e_L$ by $\Theta(\delta^{3})\cdot\Theta(\delta)=\Theta(\delta^{4})$: each
$\delta^3$-atom sits in a $D$-block whose conditional law of $C$ it shifts at relative
scale $\Theta(\delta)$ --- the atom $(0,1,1,0)$ joins the $D{=}0$ block, where $C$'s law
is $\delta$-tilted, and the atom $(0,1,0,1)$ joins the $D{=}1$ block, where it is the
sole $\delta$-scale disturbance. The two placements enter $I(C\I D)$ and $I(C\I AB)$
asymmetrically (under $(A,B)$ both atoms sit in their own block $AB{=}01$, where together
they leave $C$'s conditional law \emph{exactly} equal to $(2,27)/29$), and the
$\Theta(\delta^4)$ mismatch has a strictly negative sign for small $\delta$, which the
interval certificates of Section~\ref{sec:certificates} pin down together with its
constant.

Combining (i) and (ii): $R_L = \Theta(\delta^{4})/\Theta(\delta^{5}) =
\Theta(1/\delta)$.
\end{proof}

\begin{remark}
The certified computation below measures the sharp law: $R_L\,\delta \to 2$, i.e.\
$R_L = 2\cdot 10^{L} - c + o(1)$ with $c \approx 0.6437$; and the per-step ratio
$R_{L+1}/R_L$ is certified $\ge 1.9$ (measured: $\to 10$) at every step $2\le L\le 15$.
Theorem~\ref{thm:main} needs only $R_L\to\infty$.
\end{remark}

\section{Machine certificates}\label{sec:certificates}

The published verifier reproduces, from the integer weights alone, with directed-rounding
decimal interval arithmetic (every logarithm bracketed one-sidedly; every certified sign
a theorem about the printed digits):

\begin{center}
\begin{tabular}{rl}
\toprule
$L$ & certified lower bound on $R_L$ \\
\midrule
2 & $189.4773957$ \\
3 & $1{,}998.356506$ \\
4 & $19{,}999.25622$ \\
8 & $199{,}999{,}999.3$ \\
12 & $1.999999999\cdot 10^{12}$ \\
16 & $1.999999999\cdot 10^{16}$ \\
\bottomrule
\end{tabular}
\end{center}

\noindent together with $y_L>0$ and $e_L<0$ for $2\le L\le 16$, the per-step growth
bound, and explicit kill certificates $f_k(W_L)<0$ including
\[
k=(1,1,1)\;\text{at}\;L{=}2,\qquad
k=(10^{9},10^{9},10^{9})\;\text{at}\;L{=}15,\qquad
k=(10^{30},100,10^{30})\;\text{at}\;L{=}5,
\]
\[
k=(0,10^{20},0)\;\text{at}\;L{=}22,\qquad
k=(0,10^{24},0)\;\text{at}\;L{=}26 .
\]
By~\eqref{eq:collapse} a certificate at $k_2$ covers all $(k_1,k_2,k_3)$, so every
conjectured constant triple with $k_2\le 10^{24}$ is refuted by a named, explicit
distribution; the verifier accepts an arbitrary triple and searches out the certifying
$L$. Four independently written implementations (interval engines in
independent codebases, no shared code) agree on every row.

\section{Discussion}

What survives, and what closes. The Linden--Winter constrained inequality (their
Theorem~1) is untouched, as are the infinite constrained families of
Cadney--Linden--Winter~\cite{CLW2012}. What closes is the named route from the
constrained inequality to an unconstrained one: no linear relaxation of the constraints
works. This is consistent with, and sharpens, the current expert assessment that no
unconstrained non-Shannon-type inequality for the von Neumann entropy is in
sight~\cite{HHR2024}: the one standing candidate is now gone, through states as simple as
six-atom classical distributions on four bits.

The mechanism is reminiscent of the robustness failures of quantum Markov chains
identified by Ibinson, Linden and Winter~\cite{ILW2008}: smallness of conditional mutual
information does not control the quantities one would like it to control. Here the
failure is sharper --- the deficit $I(C\I AB)-I(C\I D)$ is not merely uncontrolled but
exceeds \emph{any} fixed linear multiple of the constraint violations.

\begin{thebibliography}{9}
\bibitem{LiebRuskai1973} E.~H. Lieb and M.~B. Ruskai,
\emph{Proof of the strong subadditivity of quantum-mechanical entropy},
J.~Math.~Phys.\ 14:1938--1941 (1973).

\bibitem{Pippenger2003} N.~Pippenger,
\emph{The inequalities of quantum information theory},
IEEE Trans.\ Inf.\ Theory 49(4):773--789 (2003).

\bibitem{LindenWinter2005} N.~Linden and A.~Winter,
\emph{A new inequality for the von Neumann entropy},
Commun.\ Math.\ Phys.\ 259:129--138 (2005). arXiv:quant-ph/0406162.

\bibitem{ILW2008} B.~Ibinson, N.~Linden and A.~Winter,
\emph{Robustness of quantum Markov chains},
Commun.\ Math.\ Phys.\ 277:289--304 (2008). arXiv:quant-ph/0611057.

\bibitem{CLW2012} J.~Cadney, N.~Linden and A.~Winter,
\emph{Infinitely many constrained inequalities for the von Neumann entropy},
IEEE Trans.\ Inf.\ Theory 58(6):3657--3663 (2012). arXiv:1107.0624.

\bibitem{HHR2024} T.~He, V.~E. Hubeny and M.~Rota,
\emph{Inner bounding the quantum entropy cone with subadditivity and subsystem coarse
grainings}, Phys.\ Rev.\ A 109:052407 (2024). arXiv:2312.04074.
\end{thebibliography}

\end{document}
